\documentclass[10pt,twocolumn]{article}
\usepackage[T1]{fontenc}
\usepackage[utf8]{inputenc}
\usepackage{lmodern}
\usepackage[margin=1.8cm,top=2.4cm,bottom=2cm,columnsep=0.6cm]{geometry}
\usepackage{fancyhdr}
\usepackage{titlesec}
\usepackage{booktabs}
\usepackage{amsmath,amssymb,amsfonts}
\usepackage{microtype}
\usepackage{xcolor}
\usepackage{mdframed}
\usepackage{parskip}
\usepackage{enumitem}
\usepackage{hyperref}

\definecolor{rulered}{RGB}{180,30,30}
\definecolor{boxgray}{RGB}{245,245,245}
\definecolor{keywordgray}{RGB}{100,100,100}

\pagestyle{fancy}
\fancyhf{}
\fancyhead[L]{\textsc{\small Claude Mag}}
\fancyhead[C]{\small\textit{Machine Learning Research Digest}}
\fancyhead[R]{\small 2026-08-03}
\fancyfoot[C]{\small\thepage}
\renewcommand{\headrulewidth}{0.8pt}

\titleformat{\section}{\normalsize\bfseries\color{rulered}}{}{0pt}{}%
  [\vspace{-4pt}\color{rulered}\rule{\columnwidth}{0.4pt}]
\titlespacing{\section}{0pt}{8pt}{4pt}

\hypersetup{colorlinks=true,urlcolor=rulered,linkcolor=black}

\begin{document}
\setlength{\parindent}{0pt}
\setlength{\parskip}{4pt}

{\small\color{keywordgray}\textbf{Keywords:}
  alignment tax \textbullet{}
  domain adaptation \textbullet{}
  parametric memory \textbullet{}
  catastrophic forgetting}\par\vspace{4pt}

{\LARGE\bfseries\raggedright
  MemSFT: Mitigating Alignment Tax with an
  External Parametric Memory}\par\vspace{4pt}

{\large\itshape\raggedright
  A Plug-and-Play Memory Module Imparts Domain Knowledge Without
  Touching Backbone Weights---Matching Full SFT While Preserving
  General Capability}\par\vspace{6pt}

{\small\textbf{By} Jiarui Wang, Xiang Shi, Jiaqi Cao, Rubin Wei,
  Xiquan Wang, Hao Sun, Jingzhi Wang, Zhiqi Yang, Qipeng Guo,
  Bowen Zhou, Zhouhan Lin
  (Shanghai Jiao Tong University; Fudan University; Tencent AI Lab)}\par\vspace{2pt}

{\small\color{keywordgray}arXiv:2607.25614 \textbullet{} July 28, 2026}
\par\vspace{2pt}\noindent{\color{rulered}\rule{\columnwidth}{0.6pt}}\vspace{6pt}

Fine-tuning an LLM on a specialized domain reliably degrades its
general-purpose capability---a phenomenon called the alignment tax.
\textbf{MemSFT} avoids this by never updating backbone parameters: a
compact parametric memory module is trained to imitate a domain retriever,
and a learned router fuses memory and backbone distributions token-by-token
at inference. Across biology, geoscience, and law, MemSFT exceeds full
SFT on domain tasks while degrading general benchmarks by only 0.2\%---
versus a 20.5-point drop for full SFT.

\begin{mdframed}[backgroundcolor=boxgray,linecolor=rulered,linewidth=1pt,
    innerleftmargin=6pt,innerrightmargin=6pt,
    innertopmargin=6pt,innerbottommargin=6pt]
\textbf{\small AT A GLANCE}\par\vspace{4pt}
\begin{description}[leftmargin=0pt,labelsep=4pt,itemsep=2pt,parsep=0pt,
    font=\small\bfseries]
  \item[Problem:] Fine-tuning on domain data causes catastrophic forgetting
    of general capabilities, limiting the practical utility of LLM adaptation.
  \item[Why it matters:] Deploying LLMs in law, medicine, and science
    requires domain expertise without sacrificing general reasoning.
  \item[Method:] Train a small parametric memory to distill a domain retriever;
    a router dynamically fuses memory and backbone at each decoding step.
  \item[Result:] 72.6 domain avg.\ (vs.\ 71.8 for full SFT);
    74.6 general avg.\ (vs.\ 54.3 for full SFT).
  \item[Evidence:] Evaluated on biology, geoscience, law with Qwen3-8B
    through Qwen3-235B-A22B; memory transfers across model sizes.
\end{description}
\end{mdframed}

\section{The Big Picture}

Large language models have transformed the landscape of domain-specific
AI: a single pretrained model can serve as the foundation for law, medicine,
finance, and science---but only with significant customization. Full
supervised fine-tuning (SFT) on domain data is the dominant adaptation
method, but it carries a well-documented cost: catastrophic forgetting.
General benchmarks drop 15--30\% relative to the base model after full SFT
on a specialized corpus.

MemSFT recasts domain adaptation as a knowledge injection problem. Instead
of modifying the backbone, it adds a compact module that holds the domain
knowledge. The backbone remains intact; the memory supplies the domain delta;
a router blends them conditionally during generation.

\section{Why Existing Approaches Fall Short}

Three adaptation paradigms exist, each with a critical limitation:

\textbf{Full SFT:} Updates all backbone parameters. Achieves strong domain
performance but causes catastrophic forgetting (general benchmarks drop
15--30\%). Once fine-tuned, the backbone cannot easily serve multiple
domains simultaneously.

\textbf{LoRA/adapters:} Updates low-rank adapter matrices ($\ll 1\%$ of
parameters). Reduces forgetting but does not eliminate it; domain gains
are weaker than full SFT.

\textbf{Retrieval-augmented generation (RAG):} Retrieves relevant documents
at inference time. Preserves general capability but adds retrieval latency
($\sim$100--500\,ms per query), requires a maintained document store, and
does not internalize patterns that allow efficient pattern completion.

MemSFT achieves full SFT domain performance with RAG-level general
capability preservation, and with only 15\% additional inference latency
over a vanilla backbone---no retrieval system required.

\section{Core Mechanism}

MemSFT consists of two components trained sequentially without backbone modification:

\textbf{Parametric memory $M_\phi$:} A compact transformer ($\sim$75M
parameters) trained to imitate a dense retriever operating over domain
documents. Given the generation prefix $h_{1:t}$, $M_\phi$ produces a
distribution over the next token that reflects domain-relevant patterns.

\textbf{Router $g_\psi$:} A lightweight gating network ($\sim$1M parameters)
that predicts, at each decoding step, what fraction of the output
distribution should come from the memory vs.\ the backbone.

\section{Technical Formulation}

The memory is trained by knowledge distillation from a retriever:
\[
  \mathcal{L}_{\text{mem}} = D_{\mathrm{KL}}\!\bigl(
    \mathrm{Retriever}(\cdot \mid h_{1:t},\, \mathcal{D}_{\mathrm{dom}})
    \;\|\;
    M_\phi(\cdot \mid h_{1:t})
  \bigr)
\]
where $\mathcal{D}_{\mathrm{dom}}$ is the domain corpus. The retriever
is a dense retrieval model that augments a frozen LM with top-$k$ documents.

The final token distribution is:
\[
  p(w_t \mid w_{<t}) =
    g_\psi(w_{<t}) \cdot M_\phi(w_t \mid w_{<t})
    + \bigl(1 - g_\psi(w_{<t})\bigr) \cdot p_\theta(w_t \mid w_{<t})
\]
where $p_\theta$ is the frozen backbone and $g_\psi(w_{<t}) \in [0,1]$.

The router is trained to maximize the likelihood of domain data under
$p(\cdot)$, with a sparsity regularizer:
\[
  \mathcal{L}_{\text{router}} =
    -\mathbb{E}_{(w,c)\sim\mathcal{D}_{\mathrm{dom}}}\!\bigl[
      \log p(w \mid c)
    \bigr] + \lambda \|g_\psi\|_1
\]

\section{Learning or Inference Procedure}

Training proceeds in two stages:

\textbf{Stage 1 (memory distillation):} Train $M_\phi$ against the
retriever output distribution on domain documents for 10K steps with Adam.
Backbone is frozen throughout.

\textbf{Stage 2 (router training):} Freeze $M_\phi$, train $g_\psi$ on a
held-out domain set for 1K steps. The sparsity regularizer encourages
the router to invoke memory selectively: on non-domain prompts, the router
converges to $g_\psi \approx 0$ (pure backbone).

At inference, the backbone and memory run in parallel; the router blends
their outputs. Per-step latency overhead is $\approx 15\%$ for Qwen3-8B.

\section{What the Guarantee Says}

Since the backbone parameters $\theta$ are never modified, the backbone's
general-task performance is preserved exactly by construction. Any
degradation in general benchmarks must arise from the router incorrectly
invoking domain memory for non-domain prompts. The sparsity regularizer
controls this: the expected general degradation is bounded above by
$\lambda \cdot \mathbb{E}[g_\psi(w_{<t})]$ on non-domain contexts, which
converges to near-zero with appropriate $\lambda$.

\section{Experimental Findings}

\textbf{Domains:} Biology (PubMedQA, BioASQ), Geoscience (GeoQA), Law
(LegalBench). Models: Qwen3-8B to Qwen3-235B-A22B.

\begin{center}\small
\begin{tabular}{lcc}
\toprule
Method & Domain avg. & General avg. \\
\midrule
Base LLM & 62.3 & 74.8 \\
Full SFT & 71.8 & 54.3 \\
LoRA & 68.4 & 70.1 \\
RAG & 70.2 & 74.7 \\
MemSFT & \textbf{72.6} & \textbf{74.6} \\
\bottomrule
\end{tabular}
\end{center}

MemSFT exceeds full SFT on domain tasks by 0.8 points while preserving
general capability (only 0.2-point drop versus a 20.5-point drop for
full SFT). It also exceeds RAG by 2.4 domain points with 3$\times$ lower
inference latency.

\textbf{Cross-model reuse:} A memory trained with Qwen3-8B transfers to
Qwen3-72B, recovering 92\% of per-domain performance vs.\ a native
Qwen3-72B memory---demonstrating that the memory is architecturally
portable across backbone sizes.

\section{Ablations and Interpretation}

\textbf{Router sparsity:} Removing $\lambda\|g_\psi\|_1$ causes the router
to always blend both distributions ($g_\psi \approx 0.5$), degrading general
performance by an additional 3.1 points while providing only marginal domain
gains.

\textbf{Memory size:} Scaling from 25M to 100M parameters improves domain
performance by 2.4 points; diminishing returns beyond 75M.

\textbf{Distillation teacher quality:} BM25 retriever (sparse) as teacher
vs.\ dense retrieval reduces domain performance by 3.8 points, confirming
that retrieval quality is the primary bottleneck for memory fidelity.

\section*{Reference}

Jiarui Wang, Xiang Shi, Jiaqi Cao, Rubin Wei, Xiquan Wang, Hao Sun,
Jingzhi Wang, Zhiqi Yang, Qipeng Guo, Bowen Zhou, Zhouhan Lin.
\textbf{MemSFT: Mitigating Alignment Tax with an External Parametric
Memory}.
arXiv:2607.25614, July 28, 2026.
\url{https://arxiv.org/abs/2607.25614}

\end{document}
